\section{最小二乘、范数与二次规划}
\label{sec:09-Convex-Optimization/04-Convex-Problems/02}

一批设备的校准数据，十二个点，理论上应该落在一条直线附近。你调了 \texttt{lstsq}，画出拟合线，然后发现它明显不对：线被抬得很高，斜率被压得很平，主群里几乎每个点都落在线的下方。翻回原始数据，两条记录在录入时被放大了好几倍。

这不是代码写错了。最小二乘按照你给它的指令，一丝不苟地把\emph{平方}残差之和压到了最小，而那两个偏出好几倍的点，平方之后贡献了目标值的一大半。求解器忠实地把线拽向它们。你以为自己在``拟合数据''，实际上你在那一行 \texttt{lstsq} 里声明了一件事：偏差越大，我越在乎，而且在乎的程度按平方增长。

\begin{center}
\begin{tikzpicture}[x=1.5cm,y=0.9cm,font=\small,>=stealth]
  \draw[->,black!55] (-0.1,0) -- (6.2,0);
  \draw[->,black!55] (0,-0.1) -- (0,4.3);
  \draw[black!45,dashed,thick] (0,1.658) -- (5.6,2.718);
  \draw[black,thick] (0,0.572) -- (5.6,3.278);
  \node[font=\footnotesize,black!55] at (6.0,2.62) {\(\ell_2\)};
  \node[font=\footnotesize] at (6.0,3.35) {\(\ell_1\)};
  \fill[black!75] (0.3,0.70) circle (1.8pt);
  \fill[black!75] (0.8,0.85) circle (1.8pt);
  \fill[black!75] (1.3,1.20) circle (1.8pt);
  \fill[black!75] (1.8,1.35) circle (1.8pt);
  \fill[black!75] (2.3,1.70) circle (1.8pt);
  \fill[black!75] (2.8,1.85) circle (1.8pt);
  \fill[black!75] (3.3,2.20) circle (1.8pt);
  \fill[black!75] (3.8,2.30) circle (1.8pt);
  \fill[black!75] (4.3,2.65) circle (1.8pt);
  \fill[black!75] (4.8,2.90) circle (1.8pt);
  \draw[black,thick] (0.6,3.50) circle (2.4pt);
  \draw[black,thick] (1.1,3.85) circle (2.4pt);
  \node[font=\footnotesize,black!70] at (2.75,3.95) {两条录错的记录};
  \draw[->,black!55] (1.95,3.92) -- (1.35,3.86);
\end{tikzpicture}

\emph{同一组点、同一族直线，只换了度量残差的方式：\(\ell_2\) 被两个离群点抬起并压平，\(\ell_1\) 几乎贴着主群走。}
\end{center}

图里两条线用的是同一个 \(A\)、同一个 \(b\)，差别只在目标函数写的是 \(\sum_i r_i^2\) 还是 \(\sum_i\abs{r_i}\)。它们都是凸问题，都有全局最优解，都能在毫秒内解完。选哪一个不是计算问题，是你对误差来源的判断。

\subsection{最小二乘是凸优化里被解得最透的问题}

先把这个起点讲清楚。超定方程 \(Ax\approx b\) 的平方拟合写成

\[
\min_x\;\norm{Ax-b}_2^2,
\]

展开是 \(x^{\trans}A^{\trans}Ax-2b^{\trans}Ax+b^{\trans}b\)，Hessian 为 \(2A^{\trans}A\succeq0\)。无论 \(A\) 长什么样，这个函数都凸，不需要任何附加假设。

\begin{proposition}
\(x^\star\) 是 \(\norm{Ax-b}_2^2\) 的全局最小点，当且仅当 \(A^{\trans}(Ax^\star-b)=0\)。
\end{proposition}

骨架两步。目标可微且凸，所以驻点即全局最优——凸性把``梯度为零''从必要条件升格成了充要条件，这一步用的是\hyperref[sec:09-Convex-Optimization/01-Optimization-Foundations/02]{``局部更好为何还不够''}里那条定理。梯度算出来是 \(2A^{\trans}(Ax-b)\)，令其为零即得。条件对账：用到了可微与凸，没有用到 \(A\) 满列秩，也没有用到解唯一；秩亏时这组\textbf{正规方程}仍然刻画了全部最优解，只不过最优解构成一个仿射子空间。

方程 \(A^{\trans}(b-Ax^\star)=0\) 的几何读法是残差与 \(A\) 的每一列都正交，也就是 \(Ax^\star\) 是列空间中离 \(b\) 最近的点，最短连线垂直于子空间，与\hyperref[sec:03-Linear-Algebra/03-Orthogonality/02]{``投影与最近点：从一条直线走向一般子空间''}讲的是同一件事。这是凸优化里少有的、能把最优解写成闭式的问题，也是后面一切复杂模型的参照系。

不过闭式解和好的数值实现是两回事。直接算 \((A^{\trans}A)^{-1}A^{\trans}b\) 会把 \(A\) 的条件数平方，\(A\) 只要略微病态，正规方程就会病得厉害。稳妥的做法是对 \(A\) 本身做 QR 分解或 SVD，绕开 \(A^{\trans}A\)，理由与实现见\hyperref[sec:08-Numerical-Analysis/05-Numerical-Linear-Algebra/02]{``QR 与最小二乘的稳定解法''}。凸建模负责说清最优点满足什么方程，数值线性代数负责可靠地走到那里，两件事分开想。

秩亏时还有一个建模缺口。若 \(A\) 的零空间非平凡，无穷多组参数给出完全相同的预测，最小二乘目标对它们一视同仁。想从中挑出范数最小的那一组，得往目标里再加一项——这就已经越出纯最小二乘，进入下一小节了。

\subsection{范数的选择是在声明你怕哪种错误}

一般的范数拟合写成 \(\min_x\norm{Ax-b}\)。范数凸，与仿射映射复合仍凸，所以这一族问题永远是凸的。但不同的范数在建模层面说的话完全不同。

\(\ell_2\) 把残差平方后相加。一个残差翻倍，代价变成四倍，因此目标函数的大部分注意力都花在最大的那几个残差上。若误差确实来自高斯噪声，这正是极大似然想要的加权方式；若误差里混着录入错误、传感器掉线或者标注失误，同样的加权方式就成了软肋——图里那两个空心点把整条线拽走了。

\(\ell_1\) 把残差取绝对值后相加。残差翻倍，代价也只翻倍，单个点能施加的拉力有上限。图里的实线因此几乎不理会那两个错误记录。代价有两处：目标在残差为零处不可微，需要次梯度或者转成线性规划来处理，见\hyperref[sec:09-Convex-Optimization/06-Optimization-Algorithms/03]{``不可微时次梯度为何仍能前进''}；同时它倾向于让一部分残差精确为零，也就是让拟合线严格穿过若干个点。这种``把一些分量顶到零''的倾向换到参数上就是稀疏性，机制见\hyperref[sec:09-Convex-Optimization/07-Applications/02]{``\(L_1\) 正则为何产生稀疏''}。

\(\ell_\infty\) 只看最大的那一个残差，其余的完全不计。它适合的场合很窄也很明确：任何单点上的偏差都不能超标，比如一条要装进硬件的近似曲线。它对离群点比 \(\ell_2\) 还敏感，因为整个目标就是那个最坏点。

上一篇的辅助变量技巧在这里立刻兑现：\(\ell_1\) 与 \(\ell_\infty\) 拟合都能写成线性规划，\(\ell_2\) 范数约束 \(\norm{Ax-b}_2\le t\) 则是标准的二阶锥约束。同一个拟合故事，换一个范数就换一类求解器，而三类都在成熟工具的覆盖范围内。所以剩下要判断的只有一件事：数据里的大偏差究竟是信号，还是噪声。

\subsection{加上约束或加上曲率，就成了二次规划}

把平方目标一般化，再允许线性约束，得到凸\textbf{二次规划}：

\[
\begin{aligned}
\text{minimize}\quad&\tfrac12x^{\trans}Px+q^{\trans}x+r\\
\text{subject to}\quad&Gx\le h,\\
&Ax=b,
\end{aligned}
\]

其中 \(P\succeq0\)。二次项能表达的东西很多：能量、方差、相邻分量之间的平滑度、偏离参考方案的代价。\(P\) 的非对角元素承载变量之间的耦合——两个变量同时增大，代价可能超过各自单独增大之和。投资组合里的协方差矩阵、控制里的状态代价、图像去噪里的差分算子，写出来都是这个形状。

\(P\succeq0\) 这个条件不能省。看到``平方项''就断定凸是常见的错误，交叉项完全可以把 Hessian 变成不定的：\(f(x_1,x_2)=x_1^2-x_2^2\) 沿一个方向向上弯、沿另一个方向向下弯，原点是鞍点而非极小点。判断凸性要看 \(P\) 的特征值，不是看式子里有没有平方。

Ridge 回归是这一类里最常用的成员：

\[
\min_w\;\tfrac12\norm{Xw-y}_2^2+\tfrac{\lambda}{2}\norm{w}_2^2 .
\]

Hessian 是 \(X^{\trans}X+\lambda I\)，只要 \(\lambda>0\) 就正定，即使 \(X\) 秩亏也给出唯一解。这一项在两个层面同时起作用：统计上它压住系数规模、限制模型复杂度；优化上它把所有平坦方向的曲率抬到至少 \(\lambda\)，于是问题变成强凸的，条件数有了上界，梯度类算法的收敛速率也随之确定。两个效果由同一个 \(\lambda I\) 提供，曲率与稳定性之间的这层关系见\hyperref[sec:09-Convex-Optimization/03-Convex-Functions/03]{``曲率区分凸、严格凸与强凸''}，统计侧的读法见\hyperref[sec:09-Convex-Optimization/07-Applications/01]{``光滑经验风险与 \(L_2\) 正则化''}。

约束版本同样常见。要求权重非负且和为一，就得到投资组合问题；要求预测值落在一个区间里，就得到带箱约束的回归。加上约束以后闭式解一般不复存在，但问题仍在标准形式之内，交给二次规划求解器即可，返回的解还附带一组乘子，可以按\hyperref[sec:09-Convex-Optimization/05-Duality/03]{``KKT 把最优性写成一组方程''}逐条核验。

\subsection{辅助变量把范数封装成一条锥约束}

范数目标 \(\min_x\norm{Ax-b}_2\) 可以写成

\[
\begin{aligned}
\text{minimize}\quad&t\\
\text{subject to}\quad&\norm{Ax-b}_2\le t.
\end{aligned}
\]

多了一个变量，目标却变成了线性的，范数这个复杂结构被整个塞进一条标准锥约束里。现代锥求解器内部做的就是这件事：把用户写的高层表达式层层展开，直到只剩线性约束、二阶锥约束和半正定锥约束，然后对这个统一形式求解。这也解释了为什么建模语言能接受那么灵活的写法——它们在编译阶段替你完成了辅助变量的引入。

辅助变量本身不构成代价，需要盯住的是三件别的事：变换后的问题与原问题是否严格等价（最优值相同，且能从新解还原出原解），变换有没有暴露出求解器能利用的结构，以及新变量的数值尺度会不会与原变量差出几个数量级。

最后提醒一句。这一篇里的问题都能解到最优，但优化意义上的最优不等于模型可靠。最小二乘的解可以检查残差是否与 \(A\) 的列正交，不正交就说明还没收敛或者实现有问题；范数拟合应该看残差的分布形态而不只是那个总和，一个巨大的残差很容易被一堆小残差在平均里抹平；正则化路径应该随 \(\lambda\) 连续变化，某处的突跳往往在提示多重最优或者数据内部的冲突。一个极小的训练误差如果来自巨大系数之间的正负抵消，它在数学上确实是最优解，在预测和数值上却经不起一点扰动。

拟合问题之所以顺手，是因为它们生下来就是凸的。下一篇要处理相反的情形：问题写出来明明不凸，却未必没有凸的写法。

